% Part: computability
% Chapter: computability-theory
% Section: rice-theorem

\documentclass[../../../include/open-logic-section]{subfiles}

\begin{document}

\olfileid{cmp}{thy}{rce}
\olsection{رائس دا قضیہ}

سوچو تاں پتا لگدا اے کہ
\olref[tot]{prop:total} دے
ثبوت وچ $\fn{Tot}$ دیاں
خاص تفصیلاں دا کوئی کردار
نہیں۔ اسی $h(x,y)$ نوں
ایداں بنایا سی کہ ایہ
مستقل فنکشن
$j(y) = 0$ وانگ کم کرے
عین اُدوں جدوں
$x$، $K$ وچ ہووے؛
پر اوس حال وچ اسی ایہنوں
کسے وی ہور جزوی قابلِ حساب
فنکشن وانگ کم کروا سکدے
ساں۔ ایس گل توں اک ہور
عام قضیہ ملدا اے۔ موٹے طور اُتے
ایہ آکھدا اے کہ قابلِ حساب فنکشناں
دی جیہڑی خاصیت کجھ فنکشناں وچ ہووے
تے کجھاں وچ نہ ہووے، اوہدا فیصلہ
قابلِ حساب طریقے نال نہیں ہو سکدا۔

یاد رکھو کہ $\cfind{0}$،
$\cfind{1}$، $\cfind{2}$،
~\dots جزوی قابلِ حساب
فنکشناں دی ساڈی معیاری
گنتی وار فہرست اے۔

\begin{thm}[رائس دا قضیہ]
  $C$ جزوی قابلِ حساب فنکشناں
  دا کوئی وی سیٹ ہووے،
  تے $A =
  \Setabs{n}{\cfind{n} \in C}$
  من لو۔ جے $A$ قابلِ حساب
  اے، تاں یا $C$ خالی اے
  یا $C$ سارے جزوی قابلِ حساب
  فنکشناں دا سیٹ اے۔
\end{thm}

{\em اشارییاں دا سیٹ} اوہ
سیٹ $A$ اے جیہدی خاصیت
ایہ اے: جے $n$ تے $m$
ایکو فنکشن «حساب» کرن
والے اشاریے ہون، تاں
یا دونوں $n$ تے $m$،
$A$ وچ ہوندے نیں، یا
کوئی وی نہیں ہوندا۔
ایہ ویکھنا اوکھا نہیں کہ قضیے وچ
دتے سیٹ~$A$ وچ ایہ خاصیت اے۔
الٹ پاسے، جے $A$~
اشارییاں دا سیٹ اے تے
$C$~اوہناں اشارییاں
نال حساب کیتے فنکشناں
دا سیٹ اے، تاں
$A = \Setabs{n}{\cfind{n} \in C}$۔

\begin{explain}
ایس اصطلاح وچ رائس دا
قضیہ ایہو آکھدا اے کہ
اشارییاں دے ایداں دے
سیٹ دا فیصلہ ممکن نہیں
جیڑا نہ خالی ہووے، نہ
سارے اشارییاں دا سیٹ۔
قضیہ سمجھݨ لئی
\emph{پروگراماں}، مثلاً
تہاڈی پسند دی پروگرامنگ
زبان وچ لکھے پروگراماں،
تے اوہناں ولوں حساب کیتے
فنکشناں وچ فرق یاد رکھنا
مدد کردا اے۔ پروگرام
(اشاریے) نحوی چیزاں نیں؛
اوہناں بارے کئی سوالاں
دا جواب قابلِ حساب اے:
کی ایس پروگرام وچ 150
توں ودھ نشان نیں؟ کی
22 توں ودھ سطراں نیں؟
کی ایس وچ «while»
بیان اے؟ کی «hello world»
لڑی کدے «print» بیان
دا آرگومنٹ بن کے آؤندی اے؟
رائس دا قضیہ آکھدا اے
کہ پروگرام دے
\emph{روَیّے} بارے
ایہو جیہے کسے وی سوال دا جواب
قابلِ حساب نہیں جیہدا جواب کجھ
پروگراماں لئی ہاں تے کجھ لئی ناں ہووے۔
مثلاً: کی پروگرام انپٹ
$0$ اُتے رُکدا اے؟
کی ایہ کدے رُکدا اے؟
کی ایہ کدے کوئی جفت
عدد آؤٹ پٹ کردا اے؟
\end{explain}

\begin{proof}[رائس دے قضیے دا ثبوت]
من لو $C$ نہ خالی اے
تے نہ سارے جزوی قابلِ
حساب فنکشناں دا سیٹ؛
$A$ اوہناں فنکشناں دے
اشارییاں دا سیٹ اے
جو~$C$ وچ نیں۔ اسی
وکھاواں گے کہ جے $A$
قابلِ حساب ہووے تاں
رُکّن دا مسئلہ حل ہو
جاوے؛ ایس لئی $A$~
قابلِ حساب نہیں۔

عامیت برقرار رکھدیاں اسی
من سکدے آں کہ اوہ فنکشن
$f$، جیڑا کسے وی انپٹ اُتے
تعریف شدہ نہیں،
$C$ وچ نہیں۔ جے
ایہ اوس سیٹ وچ ہووے تاں
ہِٹھلی دلیل وچ $C$ تے
اوس دا متمم آپس وچ
بدل دیو۔ $g$، $C$~
وچلا کوئی وی فنکشن
ہووے۔ خیال ایہ اے کہ
جے اسی~$A$ دا فیصلہ
کر سکدے، تاں $f$ نوں
حساب کرن والے اشاریے
تے $g$ نوں حساب کرن
والے اشاریے وکھ کر
سکدے؛ فیر ایس نال
رُکّن دا مسئلہ حل
کر سکدے۔

ایہو طریقہ اے۔ ہمہ گیر جزوی
قابلِ حساب فنکشناں دی مدد نال
اسی اک فنکشن ایداں تعریف کر سکدے آں:
\[
h(x,y) \simeq
\begin{cases}
\text{غیر تعریف شدہ} & \text{جے $\cfind{x}(x) \fundefined$} \\
g(y) & \text{ہور صورت وچ۔}
\end{cases}
\]
$h$ دی قدر کڈھن لئی
پہلے $\cfind{x}(x)$
حساب کرن دی کوشش
کرو؛ جے ایہ حسابی
کارروائی رُک جائے
تاں $g(y)$ حساب
کرو؛ تے جے
{\em اوہ} کارروائی
وی رُک جائے تاں
اوہدا آؤٹ پٹ دے
دیو۔ ہور رسمی طور
اُتے ایہ لکھ سکدے
آں:
\[
h(x,y) \simeq \Proj{2}{0}(g(y),\fn{Un}(x,x)).
\]
جتھے $\Proj{2}{0}(z_0, z_1) = z_0$
$2$ آرگومنٹاں والا
پروجیکشن فنکشن اے؛
ایہ $0$-واں آرگومنٹ
واپس دیندا اے تے
قابلِ حساب اے۔

تاں $h$ جزوی قابلِ
حساب فنکشناں دی
ترکیب اے؛ سجّا پاسا
اُدوں تے صرف اُدوں
تعریف شدہ اے تے~$g(y)$
دے برابر اے جدوں
$\fn{Un}(x,x)$ تے
$g(y)$ دونوں تعریف
شدہ ہون۔

دھیان دیو، مقرر~$x$
لئی جے $\cfind{x}(x)$
غیر تعریف شدہ اے
تاں $h(x,y)$ ہر~$y$
لئی غیر تعریف شدہ
اے؛ تے جے
$\cfind{x}(x)$ تعریف
شدہ اے تاں
$h(x,y) \simeq g(y)$۔
سو $x$ دی ہر مقرر
قدر لئی یا $h(x,y)$
عین $f$ وانگ کم
کردا اے یا عین $g$
وانگ۔ $\cfind{x}(x)$
تعریف شدہ اے یا نہیں،
ایہ فیصلہ کرنا ایہناں
دو حالتاں وچوں اک
دا فیصلہ کرنا اے۔
پر ایہ اوہو سوال
اے کہ $h_x(y) \simeq h(x,y)$
نال دتا فنکشن~$C$
وچ اے یا نہیں؛ جے
$A$ قابلِ حساب ہووے
تاں اسی ایہ فیصلہ
کر سکدے۔

ہور رسمی طور اُتے،
چونکہ $h$~دو آرگومنٹاں
والا جزوی قابلِ حساب
فنکشن اے، ایس لئی
ایہ $\cfind{e}$ دے
برابر اے، جتھے $e$
اوہدا اشاریہ اے۔
$s$-$m$-$n$
قضیے مطابق اک بنیادی
بازگشتی فنکشن~$s$
اے ایداں کہ ہر~$x$
لئی $\cfind{s(e,x)}(y) = h_x(y)$۔
ہن ہر $x$ لئی،
جے $\cfind{x}(x) \fdefined$
تاں $\cfind{s(e,x)}$
اوہو فنکشن اے جو~$g$،
اس لئی $s(e,x)$،
$A$ وچ اے۔ دوجے
پاسے، جے
$\cfind{x}(x)
\uparrow$ تاں
$\cfind{s(e,x)}$
اوہو فنکشن اے جو~$f$،
اس لئی $s(e,x)$،
~$A$ وچ نہیں۔
ہور لفظاں وچ،
ہر~$x$ لئی،
$x \in K$ جے تے
صرف جے $s(e,x) \in A$۔
جے $A$ قابلِ حساب
ہووے تاں $K$~وی
قابلِ حساب ہووے گا،
جو تضاد اے۔ ایس
لئی $A$ قابلِ حساب
نہیں۔
\end{proof}

رائس دا قضیہ بہت
طاقتور اے۔ ہِٹھلا
فوری نتیجہ ایس دیاں
کجھ مثالاں دیندا اے۔
\begin{cor}
ہِٹھاں دتے سیٹ
ناقابلِ فیصلہ نیں۔
\begin{enumerate}
\item $\Setabs{x}{\text{$17$، $\cfind{x}$ دے حاصل قدراں وچ اے}}$
\item $\Setabs{x}{\text{$\cfind{x}$ مستقل اے}}$
\item $\Setabs{x}{\text{$\cfind{x}$ کُل اے}}$
\item $\Setabs{x}{\text{جدوں وی $y < y'$، $\cfind{x}(y) \fdefined$ ہووے،
    تے جے $\cfind{x}(y') \fdefined$ ہووے تاں $\cfind{x}(y) < \cfind{x}(y')$}}$
\end{enumerate}
\end{cor}

\begin{proof}
  ایہ سارے اشارییاں
  دے اوہ سیٹ نیں جیڑے
  نہ خالی نیں، نہ سارے
  اشارییاں اُتے مشتمل۔
\end{proof}

\end{document}
